% Restructure per advisors (2026-09-11): Intro positioning + contributions, new §II Goal and Challenges,
% new §III FISH (challenges answered; expansion operator + GPU DAG live here). §IV/§V unchanged.
% Figure numbering assumed: Fig.~\ref{fig:model} = the model figure (now first), Fig.~\ref{fig:viz} = FISH-eye.

%%%%%%%%%%%%%%%%%%%%%%%%%%%%%%%%%%%%%%%%%%%%%%%%%%%%%%%%%%%%%%%%%%%%%%%%%%%%%%
% §I — replace the gap paragraph + contributions
%%%%%%%%%%%%%%%%%%%%%%%%%%%%%%%%%%%%%%%%%%%%%%%%%%%%%%%%%%%%%%%%%%%%%%%%%%%%%%
In this work, we fill that gap with FISH\footnote{GitHub: \url{https://github.com/FatihTasyaran/FISH_INTERFERE}}\textsuperscript{,}\footnote{Documentation: \url{https://fish-interfere.readthedocs.io}}, a toolset that automatically extracts this information from a traced run of a ROS~2 application: it does not modify the application, applies to any application, and yields both the application's structure and its timing behaviour as deployed. The model is based on actual execution traces. Therefore, it does not require expert knowledge about the application that needs to be modeled, nor does it rely on potentially incomplete or outdated documentation. \textbf{Our contributions are:} \textbf{(I)}~a toolset that extracts, transparently and automatically, the structural and timing information a timing analysis of a deployed ROS~2 application needs, across the OS, ROS~2 and GPU layers (Sections~\ref{sec:goal} and~\ref{sec:method}); \textbf{(II)}~the instrumentation and acquisition that make this transparent: 28 tracepoints added to \texttt{ros2\_tracing} for object lifetimes, all five kinds of callback, intra-process delivery and executor configuration, for rclcpp and rclpy alike, and a launch procedure that brings GPU processes under the profiler; \textbf{(III)}~evidence of feasibility on Autoware, an industrial-size application: the extracted model and the cost of tracing (Section~\ref{sec:findings}).

%%%%%%%%%%%%%%%%%%%%%%%%%%%%%%%%%%%%%%%%%%%%%%%%%%%%%%%%%%%%%%%%%%%%%%%%%%%%%%
\section{Goal and Challenges}
\label{sec:goal}
To analyse, verify or validate the timing of a ROS~2 application, one needs information about the deployed system that is rarely written down and, when it is, rarely current. The goal of FISH is to produce that information from a traced run of the application without modifying it, so that any application can be traced as deployed. Its output is the model of Section~\ref{sec:intro}: the application's structure and its timing behaviour as it actually ran. Three kinds of information are needed, and each is hard to obtain for its own reasons.

\textbf{Architecture---what is there.} Which components exist, which callbacks each runs, what triggers them, which channels they use, and how the executors are configured: type, thread count, callback groups. Launch files and source code do not settle this. Composable nodes are loaded into containers at run time, and objects are created and destroyed while the application runs, so handles are reused (\textbf{C1}). Channel use is not declared: any callback of a node may publish on its publishers, deliveries between nodes of one process never reach the client library, and of the five kinds of callback the executor dispatches, waitables---intra-process subscriptions, action servers---are covered by no existing tool (\textbf{C2}). Executor configuration is recorded by none of them (\textbf{C3}).

\textbf{Execution---when and where each callback runs.} The analyses of multi-threaded executors need the thread a callback ran on, and so does any attribution of GPU work to the callback that submitted it; under a multi-threaded executor that thread is decided at run time (\textbf{C4}). CPU and GPU are two tracing domains with their own tracers and clocks; a callback's GPU work is known only once the two are aligned and each GPU record is tied to the callback that issued it (\textbf{C5}). Tracing must also stay transparent: the GPU profiler sees only processes started under it, and a container restarted later would come up without its composable nodes (\textbf{C6}).

\textbf{Task chains and timing---how it behaves.} Timing must be taken in the right phase: initialisation, warm-up and steady state behave differently, and not every registered callback runs at all (\textbf{C7}). A callback that submits GPU work does not do the same thing every time; its GPU work branches from invocation to invocation (\textbf{C8}). Chains have to be identified in an entangled graph: which callbacks belong to one functionality, where inputs are joined, where a callback reads state that another wrote instead of receiving a message (\textbf{C9}). Such a callback is not one segment but several---CPU work, submissions, waits for the GPU (\textbf{C10}). Finally, a model that is observed rather than assumed is only as good as its coverage, so the extraction itself must be checked (\textbf{C11}).

%%%%%%%%%%%%%%%%%%%%%%%%%%%%%%%%%%%%%%%%%%%%%%%%%%%%%%%%%%%%%%%%%%%%%%%%%%%%%%
\section{FISH}
\label{sec:method}
FISH currently targets ROS~2 Humble and has four parts: \emph{instrumentation}, tracepoints in the ROS~2 client libraries plus the CUDA profiler; \emph{acquisition}, a library and a daemon that run an unmodified application under that instrumentation and delimit the steady-state part of the trace; \emph{extraction}, which turns the trace into the model; and FISH-eye, a viewer that lets one revisit a traced run at any granularity (Fig.~\ref{fig:viz}). The parts answer the challenges of Section~\ref{sec:goal} as follows.

\textbf{Instrumentation (C1--C5):} FISH builds on \texttt{ros2\_tracing}~\cite{bedard-2022-ros2tracing}, whose initialisation events record each object as it is created and whose execution events record what then happens to it, linked by the object's handle. To the 14 stock events of Humble, FISH adds 28. Finalize events close the lifetimes that the stock init events open and a timer event binds a timer to its node and period, so a callback is attributed correctly even after a handle is reused (C1). Start and end events for waitables and clients, request-sent and response-received events that pair a client's call with its response, and action-server events complete the five kinds of callback the executor dispatches; link events at the publish, receive and client-request sites tie each message to the callback that issued it; the intra-process subscription waitable covers deliveries that never reach \texttt{rcl}; and registration events give rclpy nodes the same callback chain as rclcpp (C2). Introspection events record each executor's type and thread count and each callback's group (C3). Every execution event carries the thread it ran on, and a thread-local active-callback context attributes the GPU API calls made on that thread to the callback in progress (C4). GPU work is recorded by CUPTI through \texttt{nsys}, aligned to the LTTng clock at ingest and bound through correlation ids to the API call and thus to the callback (C5).

\textbf{Acquisition (C6, C7):} The application is not modified: the instrumented libraries are installed on the host or in a container, and the FISH entry point wraps the \texttt{ros2 launch} command. Because CUPTI sees only processes started under the profiler, the entry point reads the launch description and starts the containers that will host GPU components under \texttt{nsys} from the outset; standalone GPU processes are caught by the FISH daemon, which polls \texttt{/proc} and restarts any process holding a CUDA context under the profiler (C6). Once the node list has been unchanged over five polls and every component is loaded, the daemon starts the input: a rosbag replay, or a schedule of timed commands for a live system. Phase boundaries come from the trace: initialisation ends when every periodic timer has fired once, warm-up when every recurring callback has completed its first execution; callbacks that never run stay in the model, marked as such (C7).

\textbf{Extraction (C8--C11):} The model is a hierarchical graph (Fig.~\ref{fig:model}a): host, executors, nodes, entities---the timers, subscriptions, services, clients and waitables an executor waits on---and callbacks, each layer partitioning the one below. Initialisation events create the vertices and the edges the wiring declares: a \emph{msg} edge wherever a publisher and a subscription share a topic, a \emph{dep} pair wherever a client and a server share a service. Execution events find the callback that actually issued each interaction and add the three kinds only running reveals: state polled through \texttt{take}, state inferred from a store and a later read, and joins that complete a set of inputs (Table~\ref{tab:edges}, Fig.~\ref{fig:model}c). Every edge is lifted along the containment, so two callbacks exchanging a message are also two entities, two nodes and two executors exchanging it (Fig.~\ref{fig:model}b), and the expansion operator opens chosen vertices into their children and redistributes their edges without creating or losing an interaction: this is how the entangled graph is read at the granularity a chain needs, one executor opened into its nodes while the rest stay whole (C9). A callback that submits GPU work owns one typed DAG per invocation---CPU segments, kernels and copies as nodes; launch, CPU order and stream FIFO as edges (C10). Invocations with the same node sequence form a shape, and folding the steady-state shapes onto their longest common subsequence gives one conditional graph whose branches carry the share of invocations taking them (Fig.~\ref{fig:model}d) (C8). Finally, the graph is checked against what the traced binaries create---a node constructor must yield one node, seven entities and seven callbacks---and is complete only if every such expectation appears in it (C11).
